% Parcelación · Filosofía 10.° · Periodo IV · PILOTO (ejemplo completo del kit)
% Modo de trabajo del docente (curriculo/MODO-DE-TRABAJO.md): sin dictado (explica; apuntes firmados),
% un entregable por clase con sello, 1 a 3 talleres (máx. 3), 1 quiz, evaluación de periodo al final y
% actividades fuera de lo común. La plantilla avisa al compilar si algo no se cumple.
% 8 semanas × 2 clases = 16 clases. Sin fechas (reutilizable): \fechainicio y \diasclase{lun,jue}.
% Los códigos [D1]… remiten a los desempeños de malla-P4.tex (FIL-10-P4-D1…).
% Compilar: bash curriculo/plantilla/compilar.sh curriculo/filosofia/10/parcelacion-P4.tex
\documentclass[parcelacion]{planeacion}
\materia{filosofia}
\grado{10.°} \periodo{IV} \anio{2027}
\clasesporsemana{2}
\carpeta{}
\begin{document}
\begin{parcelacion}
  % ---- Semana 1 · D1 La herencia griega
  \clase{\textbf{Una larga herencia.} Pregunta del periodo; Platón y Aristóteles.}
    {Explicación con \emph{La escuela de Atenas} proyectada: «¿quién es quién?» (los estudiantes
     buscan a Platón, Aristóteles, Pitágoras, Diógenes…).}
    {Reproducción en alta resolución de la obra de Rafael.}
    {\firma{Apuntes de la explicación.} {\scriptsize[D1]}}
  \clase{\textbf{El viaje de los libros.} Neoplatónicos, traducciones del árabe y del griego,
    Toledo, Averroes.}
    {Retomamos el cuadro de Rafael. Explicación breve con mapa; reto en parejas (15 min):
     «traducir» una frase de Aristóteles pasando por tres lenguas con un glosario dado y ver qué se
     pierde en el camino.}
    {Mapa proyectado; glosario de una página.}
    {\sello{Mapa del viaje de los textos con fechas.} {\scriptsize[D1]}}
  % ---- Semana 2 · D2 San Agustín
  \clase{\textbf{San Agustín: la interioridad.} \emph{Confesiones} I, 1; «vuelve a ti mismo».}
    {Explicación y lectura guiada de fragmentos; los estudiantes toman apuntes.}
    {Fragmentos de Agustín con referencia.}
    {\firma{Apuntes y tesis de cada fragmento.} {\scriptsize[D2]}}
  \clase{\textbf{Una carta a Agustín.}}
    {Cada estudiante le escribe a Agustín desde 2027: ¿la verdad se busca dentro o fuera de uno?
     Lectura en voz alta de algunas cartas.}
    {Fragmentos de la clase anterior.}
    {\sello{Carta a Agustín (media página).} {\scriptsize[D2]}}
  % ---- Semana 3 · D2 → D3
  \clase{\textbf{Fe y razón.} «Comprende para creer, cree para comprender».}
    {Experimento de confianza: una caja cerrada con un objeto; ¿qué creemos antes de comprobar y
     por qué? Conversatorio sobre fe y razón.}
    {Caja cerrada; ilusiones ópticas proyectadas.}
    {\sello{Reflexión de cinco líneas: ¿creer o comprender primero?} {\scriptsize[D2]}}
  \clase{\textbf{Tomás de Aquino y la universidad medieval.} La recuperación de Aristóteles.}
    {Explicación: cómo era estudiar en París en 1260 (horarios, libros, exámenes).}
    {Presentación; imagen de una clase universitaria medieval.}
    {\firma{Apuntes de la explicación.} {\scriptsize[D3]}}
  % ---- Semana 4 · D3 La quaestio
  \clase{\textbf{Disputa medieval en vivo.} La estructura del artículo de la \emph{Suma teológica}.}
    {Recreación de una \emph{disputatio}: un «maestro» y dos «bachilleres» discuten una pregunta;
     el curso identifica objeciones, «en sentido contrario», respuesta y solución.}
    {Artículo de la \emph{Suma teológica} (fragmento con referencia); campana para los turnos.}
    {\sello{Esquema de las partes de la disputa.} {\scriptsize[D3]}}
  \clase[t]{\textbf{Taller 1: mi \emph{quaestio}.}}
    {En parejas, una \emph{quaestio} sobre un problema actual (p.~ej. «¿Es lícito usar
     inteligencia artificial en las tareas?»).}
    {Guía del taller con el modelo y los criterios.}
    {Taller evaluativo. {\scriptsize[D3]}}
  % ---- Semana 5 · D4 ¿Edad Oscura?
  \clase{\textbf{Juicio a la Edad Media (preparación).} Origen de la etiqueta «Edad Oscura».}
    {Reparto de roles: fiscalía, defensa, jurado y testigos (un monje copista, un médico de la
     peste, un traductor de Toledo, un constructor de catedrales); cada uno reúne evidencias.}
    {Fichas de personaje y de evidencias.}
    {\sello{Ficha del personaje con sus evidencias.} {\scriptsize[D4]}}
  \clase{\textbf{Juicio a la Edad Media.}}
    {Audiencia: alegatos, testigos y veredicto del jurado; cierre: ¿es justa la etiqueta?}
    {Fichas; reglas de la audiencia.}
    {\sello{Veredicto razonado (individual).} {\scriptsize[D4]}}
  % ---- Semana 6 · D4 La imprenta
  \clase{\textbf{La imprenta y la difusión de las ideas.}}
    {Carrera copista contra imprenta: cinco minutos copiando a mano un párrafo frente a una hoja
     fotocopiada; calculamos cuánto tardaría copiar un libro entero. Lectura de una gráfica sobre
     la expansión de la imprenta.}
    {Párrafo impreso; gráfica proyectada.}
    {\sello{Cálculo de la carrera y dos conclusiones sobre la gráfica.} {\scriptsize[D4]}}
  \clase[q]{\textbf{Quiz del periodo.} Agustín, Tomás y la Edad Media.}
    {Quiz y retroalimentación inmediata.}
    {Quiz impreso.}
    {Quiz. {\scriptsize[D1--D4]}}
  % ---- Semana 7 · D4 → D5
  \clase{\textbf{Humanismo: Erasmo de Rotterdam.} \emph{Elogio de la locura} (1511).}
    {«La Locura toma la palabra»: lectura dramatizada de un fragmento; en grupos, la Locura opina
     sobre algo de hoy (las redes, el colegio, la moda).}
    {Fragmento de Erasmo.}
    {\sello{Discurso breve de la Locura.} {\scriptsize[D4]}}
  \clase{\textbf{Tomás Moro y la \emph{Utopía} (1516).} Las ovejas que devoran hombres.}
    {Explicación y lectura de fragmentos del libro I y del libro II: qué critica Moro y qué
     propone.}
    {Fragmentos de la \emph{Utopía}.}
    {\firma{Apuntes y crítica central de Moro.} {\scriptsize[D5]}}
  % ---- Semana 8 · Cierre
  \clase[t]{\textbf{Taller 2: diseña tu isla de Utopía.}}
    {En grupos: mapa de la isla, tres leyes con sus razones y un riesgo de cada ley; exposición
     rápida.}
    {Guía del taller; papel y marcadores.}
    {Taller evaluativo. {\scriptsize[D5]}}
  \clase[e]{\textbf{Evaluación de periodo.} De la escolástica a la \emph{Utopía}.}
    {Evaluación escrita tipo ICFES (selección y preguntas abiertas).}
    {Evaluación impresa.}
    {Evaluación de periodo. {\scriptsize[D1--D5]}}
\end{parcelacion}
\end{document}
